%  LaTeX support: latex@mdpi.com 
%  For support, please attach all files needed for compiling as well as the log file, and specify your operating system, LaTeX version, and LaTeX editor.

%=================================================================
\documentclass[jimaging,article,submit,pdftex,moreauthors]{Definitions/mdpi} 

%=================================================================

% MDPI internal commands
\firstpage{1} 
\makeatletter 
\setcounter{page}{\@firstpage} 
\makeatother
\pubvolume{1}
\issuenum{1}
\articlenumber{0}
\pubyear{2023}
\copyrightyear{2023}
%\externaleditor{Academic Editor: Firstname Lastname}
\datereceived{23 February 2023} 
\daterevised{2 April 2023} 
\dateaccepted{} 
\datepublished{} 
\hreflink{https://doi.org/} % If needed use \linebreak
%\doinum{}

%=================================================================
% Add packages and commands here. The following packages are loaded in our class file: fontenc, inputenc, calc, indentfirst, fancyhdr, graphicx, epstopdf, lastpage, ifthen, lineno, float, amsmath, setspace, enumitem, mathpazo, booktabs, titlesec, etoolbox, tabto, xcolor, soul, multirow, microtype, tikz, totcount, changepage, attrib, upgreek, cleveref, amsthm, hyphenat, natbib, hyperref, footmisc, url, geometry, newfloat, caption

\graphicspath{{figures/}} % path to folder with figures
\usepackage{subcaption}
\usepackage{listings}
\usepackage{xcolor}
\usepackage{gensymb} % for \textdegree
\usepackage{tabularx}
\usepackage{array}
\usepackage{makecell}
\usepackage{multirow}

\usepackage{listings}
\usepackage{xcolor}

\definecolor{deepblue}{rgb}{0,0,0.5}
\definecolor{deepred}{rgb}{0.6,0,0}
\definecolor{deepmagenta}{RGB}{195,12,214}
\definecolor{deepgreen}{RGB}{29,122,30}
\definecolor{mintgreen}{RGB}{192,249,192}
\definecolor{codepurple}{RGB}{204, 0, 153}
\definecolor{LightCyan}{rgb}{0.88,1,1}
\definecolor{GainsboroPurple}{RGB}{247,176,247}
\definecolor{LightSlateGrey}{RGB}{124,118,118}
%    
\lstdefinestyle{mystyle}{
    commentstyle=\color{LightSlateGrey},
    keywordstyle=\color{deepgreen}\bfseries,
    emphstyle=\color{deepred},
    numberstyle=\tiny\color{deepmagenta},
    stringstyle=\color{codepurple},
    basicstyle=\ttfamily\scriptsize,
    breakatwhitespace=false,         
    breaklines=true,                 
    captionpos=b,                    
    keepspaces=true,                 
    numbers=left,                    
    numbersep=5pt,                  
    showspaces=false,                
    showstringspaces=false,
    showtabs=false,                  
    tabsize=2
}
\lstset{style=mystyle}

%================================================================= in Khartoum Region of Sudan, NE Africa
% Full title of the paper (Capitalized)
\Title{Multispectral Satellite Image Analysis for Computing Vegetation Indices by R in the Khartoum Region of Sudan, Northeast Africa}

% MDPI internal command: Title for citation in the left column
\TitleCitation{Multispectral Satellite Image Analysis for Computing Vegetation Indices by R in the Khartoum Region of Sudan, Northeast Africa}

% Author Orchid ID: enter ID or remove command
\newcommand{\orcidauthorA}{0000-0002-5759-1089} % Polina Add \orcidA{} behind the author's name
\newcommand{\orcidauthorB}{0000-0002-6461-1551} % Olivier Add \orcidB{} behind the author's name

% Authors, for the paper (add full first names)
\Author{Polina Lemenkova $^{1}$*\orcidA{} and Olivier Debeir $^{1}$\orcidB{}}

%\longauthorlist{yes}

% MDPI internal command: Authors, for metadata in PDF
\AuthorNames{Polina Lemenkova and Olivier Debeir}

% MDPI internal command: Authors, for citation in the left column
\AuthorCitation{Lemenkova, P.; Debeir, O.}
% If this is a Chicago style journal: Lastname, Firstname, Firstname Lastname, and Firstname Lastname.

% Affiliations / Addresses (Add [1] after \address if there is only one affiliation.)
\address{%
$^{1}$ \quad Laboratory of Image Synthesis and Analysis (LISA), École polytechnique de Bruxelles (Brussels Faculty of Engineering), Université Libre de Bruxelles (ULB). Building L, Campus du Solbosch, ULB -- LISA CP165/57, Avenue Franklin D. Roosevelt 50, Brussels 1050, Belgium.
}

% Contact information of the corresponding author
\corres{Correspondence: polina.lemenkova@ulb.be; Tel.: +32 471 86 04 59 (P.L.)}

% Current address and/or shared authorship
%\firstnote{Current address: Affiliation 3.} 

\abstract{Desertification is one of the most destructive climate-related issues in Sudan-Sahel region of Africa. The assessment of desertification is possible by satellite image analysis using Vegetation Indices (VIs). \textcolor{blue}{This study reports technical advantages and features of R scripts and packages ‘raster’ and ‘terra’ for computing the VIs. The test area includes the region of the confluence of the Blue Nile and White Nile in Khartoum, southern Sudan, Northeast Africa. The Landsat 8-9 OLI/TIRS images taken for the periods of 2013, 2018 and 2022 are chosen as test datasets.} The VIs are \textcolor{blue}{used} as robust indicators of plant \textcolor{blue}{greenness as} parameters of vegetation coverage essential for environmental \textcolor{blue}{analytics.} \textcolor{blue}{Five VIs were calculated }to compare the status and dynamics \textcolor{blue}{of} vegetation \textcolor{blue}{through the} difference \textcolor{blue}{between the images} collected \textcolor{blue}{within the} nine year time gap. \textcolor{blue}{Using scripts} for \textcolor{blue}{the} computing and visualization of \textcolor{blue}{ the }VIs \textcolor{blue}{demonstrates previously unreported patterns of vegetation over Sudan to reveal climate--vegetation relationships. Technical} advantages of \textcolor{blue}{the} R packages 'raster' and 'terra' \textcolor{blue}{are highlighted} for spatial data processing to automate image analysis and mapping \textcolor{blue}{with a case study of Sudan which presents new perspectives for image processing}.
}
% Keywords
\keyword{Remote sensing; Sudan; Khartoum; Nile River; image analysis; image processing; information retrieval; programming; Africa; computer vision} 

\PACS{91.10.Da; 91.10.Jf; 91.10.Sp; 91.10.Xa; 96.25.Vt; 91.10.Fc; 95.40.+s; 95.75.Qr; 95.75.Rs; 42.68.Wt}
\MSC{86A30; 86-08; 86A99; 86A04}
\JEL{Y91; Q20; Q24; Q23; Q3; Q01; R11; O44; O13; Q5; Q51; Q55; N57; C6; C61}

%%%%%%%%%%%%%%%%%%%%%%%%%%%%%%%%%%%%%%%%%%
\begin{document}

\section{Introduction}

\subsection{Background}

In this paper, we introduce the application of R libraries for remote sensing data processing \hl{in the context of computing vegetation indices (VIs) that indicate desertification processes. Further, we} focus on \hl{calculating }several satellite-derived VIs \hl{and discuss} their technical details and differences\hl{. The test dataset include the} Landsat 8-9 OLI/TIRS images collected on Khartoum city around the confluence of the White Nile and the Blue Nile.\hl{ }The \hl{images were used to }analyse\hl{ }the changes in \hl{the }vegetation patterns and greenness \hl{through the }comparison of the time series and \hl{visualisation of} the dynamics in vegetation coverage. The standard\hl{ methods }for \hl{such tasks} are based on using the traditional \hl{GIS }software\hl{ }\cite{AHMED201829} and involve many manual \hl{routine }operations \cite{Hawash,YOUSSEF2020103777,6621869,Aldoma}. Instead of using such conventional approaches, we \hl{apply} several libraries of R such as 'raster' and 'terra' \hl{and} auxiliary packages in a principled way to obtain a better workflow for image \hl{analysis.} 

\hl{The Earth observing satellite sensors provide the remote sensing data useful in environmental studies. }The \hl{benefits} of the use of the remote sensing data and image processing are discussed earlier \cite{su14137871,ELBASTAWESY2017252,SOSNOWSKI201651,GORSEVSKI201364}. \hl{Satellite images can be used to analyse the links between the complex hydrological and climate processes and vegetation responses that lead to desertification}. For example\hl{, }Landsat images \hl{enable to }detect\hl{ }land cover types \hl{through the time-series analysis}\cite{WILSON2002385,Lemenkova202227,4293065,SALIH2017S21}, \hl{compute} vegetation indices \cite{5416935,Lemenkova202229,1026428} and \hl{show the advance or retreat of arid areas for analysis of the} desertification \cite{RIVERAMARIN2022104829}. Specifically, for the Sudan and Nile River Basin, \hl{the Landsat data are the precious source of information, due to fieldwork }data scarcity \cite{Ali} such as regular measurements of rainfall, streams runoff, and  weather stations data. This highlights the importance of the \hl{remote sensing} data for climate-related studies \hl{of Sudan.}

\hl{The GIS was one of the earliest computer-based application for satellite image processing developed for the empirical investigations of the environmental parameters derived from the remote sensing data. Since the 1990s researchers were able to accurately compute the vegetation indices derived from the images using GIS. Kumar et al. have studied performance of the ERDAS Imagine in computing vegetation indices to assess desertification in the semi-arid regions, emphasising that the NDVI classification is suitable for assessing land degradation using the NDVI pixel range} \cite{KUMAR2022100578}. \hl{Rashid et al. have examined the links between the NDVI and the salinity characteristics in the surface water, groundwater, and soils using the ArcGIS} \cite{RASHID2023100314}. \hl{Further, Ezaidi et al. applied ENVI GIS for computing the NDVI for a time series of the Landsat images to assess the degradation in vegetation} \cite{EZAIDI2022100800} \hl{These studies illustrate the traditional approaches for computing vegetation indices by GIS.}

\hl{As a general overview, GIS enables the calculation of vegetation indices from Landsat images using map algebra. However, the efficacy of the computational approaches as more advanced tools of spatial data analysis enabled to apply modelling methods for computation of the VIs. For instance, this demonstrated by Generalised Additive Models used to study the relationships between vegetation types and NDVI} \cite{MARTINEZ2023115155}, \hl{or Holt-Winters modelling to detect trend and seasonal vegetation pattern using Python's Scikit-Learn library} \cite{OMAR2021e01020}\hl{, or partial least squared regression (PLSR) model used by SIMCA-P software to track the non-linear relationships between the NDVI and climatic variables} \cite{WANG2022102704}. \hl{Recent development of the machine learning tools and Python programming enabled to use the Convolutional Neural Network (CNN) and Deep Neural Networks (DNN) for environmental analysis and estimation of the VIs} \cite{BARRIGUINHA2022103069,ROY2021100582,LI2022102818,DAVIDSON2022107396}. 

\hl{Compared to these and similar programming approaches, the use of R in the context of the computing of VIs presents a more straightforward approach. Similar to GIS approaches, R is adjusted to the processing of the Landsat bands; the same time, it benefits from the flexibility of programming methods that enables a full control on the data processing and high level of automation. The examples of R used for computing VIs present the 'raster' package for computing the NDVI }\cite{LEVEAU2021127199}\hl{, 'npphen' R package to assess the NDVI anomalies and evaluating spatio-temporal changes} \cite{CHAVEZ2023103138}\hl{, 'akima' R package for interpolation of environmental parameters} \cite{KLIMAVICIUS2023171} \hl{or the 'RStoolbox' package for computing the NDVI} \cite{PRAVALIE2022108629}. \hl{Such methods are consonant with our approach, however, they are largely tied to specific problem formulation applied to the target areas. To the best of our knowledge, no previous studies were reported on the use of R for computing the VIs in Sudan. This study fills in this gap and highlights the benefits of R packages as advanced technical tools of computing the VIs specifically for Sahelian landscapes of southern Sudan, prone to desertification.}

\subsection{Study Area}

Desertification is one of the most destructive climate-related \hl{issue} in Sudan \hl{which is located }southwards of Sahara Desert in the Sudan-Sahel region of Africa, \hl{(}Figure \ref{fig01}). \hl{The desertification} results from various factors including climate change and human activities. The \hl{climate }drivers\hl{ }include \hl{rising temperatures, decreased relative humidity} \cite{Alvi}{, increased }soil salinization \cite{DAWELBAIT201245}, seasonality \hl{and decline }of rainfall \cite{MOHAMED20223265}, \hl{instability} and irregularity of wet and dry seasons \cite{HULME198689}, the variability of precipitation \cite{Elagib,Ayoub} and the overall water deficit \cite{akhtar1994methods}. Recent studies reported \cite{SALIH2017S31} frequent drought periods during the last \hl{decades}. Severe droughts\hl{ }result in several cycles of famine in Sudan with affected farmers and population. \hl{For instance, }chronic drought \hl{is reported }in Sahel region of Darfur \cite{FULLER198739}, droughts in central Sudan \cite{Salih2022}, or eastern Sudan \cite{TARIKU2021144863,SULIEMAN2012132}. \hl{Other factors include changed patterns of cloudiness and atmospheric radiation which }results in high evapotranspiration \hl{in} the extreme north near Sahara and \hl{southern} Sahel and savannah \hl{and affects vegetation growth }\cite{ElKarouri1986}. 

These processes trigger environmental degradation\hl{, }especially in vulnerable zones of central and southern Sahelian zone of Sudan \cite{Pearce,YANG2022106213} and savannah \cite{Falkenmark}. \hl{The} most endangered region \hl{of Sudan prone to the risks of desertification lies between} 12° to 18°N, \hl{near its} capital -- Khartoum. \hl{Flat} topography of the country exaggerates the desertification problems \hl{and} increases the effects from dust and sand storms \hl{that are }often in the northern regions of \hl{Sudan} \cite{khiry2006mapping,9635812,7948717,7381366}. Specifically, this includes the effects from the \hl{creeping} sands and dunes from the areas of Sahara, Nubian and Libyan deserts, Figure \ref{fig01}. As a consequence, this results in deforestation, soil desiccation, changed \hl{land} cover types \cite{Biro} and decreased biodiversity. The latter includes, \hl{the} extinction of wild species in \hl{oases. Consequences} of land cover changes \hl{include }the imbalance of ecosystems\hl{, }spread of \hl{the }invasive species \cite{insects10110405} and disturbance of natural habitat patterns \cite{Fuller}. 

\begin{figure}[H]
\centering
	\includegraphics[width=14.0 cm]{Fig_01.jpg}
	\caption{Topographic map of Sudan. Mapping software: Generic Mapping Tools (GMT) scripting toolset version 6.1.1, \cite{Wessel}. Data source: GEBCO/SRTM. Cartography source: authors.
	\label{fig01}}
\end{figure}

70\% of the territory of Sudan belongs to the Nile basin \cite{Melesse} with 5 from the total 6 Cataracts of Nile located within Sudan (Figure \ref{fig01}). Moreover, roughly 2/3 of the territory of Nile lies within Sudan with its major tributaries -- the While Nile and Blue Nile -- joining \hl{near }Khartoum \cite{Hamad}. Hence, the specifics of the Khartoum region \hl{is} connected to the Nile River system and the confluence of the White Nile and Blue Nile, Figure \ref{fig02}. The \hl{impact} the Nile River \hl{on} the environment \hl{includes} occasional floods, soil erosion and succession of vegetation on the islands and banks \cite{Halwagy}\hl{,} accumulation of sedimentation \cite{BILLI201012}, growth of aquatic \hl{weeds} \cite{PACINI2016242}. The Nile also provides a specific \hl{place} for wetland habitat \hl{of} rare species. \hl{As a central water artery of the country, the Nile River plays a} crucial role \hl{for }hydropower \hl{of} Sudan \cite{AYIK20232229} \hl{as an important transport waterway and providing} water resources for 85\% \hl{of} population and irrigated agriculture \cite{Kauetal2023}. 

\begin{figure}[H]
\makebox[1.0\linewidth][r]{%
	\begin{subfigure}[b]{.5\textwidth}
		\centering
			\includegraphics[width=7cm]{Fig_02a.jpg}
		\caption{}
	\end{subfigure}%
	\begin{subfigure}[b]{.5\textwidth}
		\centering
			\includegraphics[width=7cm]{Fig_02b.jpg}
		\caption{}
	\end{subfigure}%
}
\caption{The confluence of the White Nile and Blue Nile in Khartoum, visualized based on the map and aerial images from the Google Earth: (\textbf{a}) Google Earth Map, (\textbf{b}) Google Earth Image.}\label{fig02}
\end{figure}

Social factors of desertification include \hl{processes} related to land management \hl{and} agriculture policymaking \cite{KOBAYASHI2020257}. Although agriculture and livestock \hl{play} an important role in the economy of Sudan \hl{as} major sources of \hl{livelihood} \cite{sulieman2010expansion}, \hl{the }unsustainable agricultural expansion and farming \hl{contribute} to \hl{desertification} \cite{Ibrahim} \hl{since} uncontrolled agriculture results in soil erosion, compaction and nutrient depletion, changed irrigation \hl{patterns} \cite{AHMED2020106064} and the increase in pesticides \cite{Nesser} \hl{which }contribute to land cover changes. Moreover, \hl{cattle} overgrazing modify \hl{landscape} structure \hl{and contribute to desertification}\cite{Ahmed,Sulieman,Akhtar}. \hl{The }intense practice of urban agriculture in Khartoum aims at meeting the increasing food demand of the \hl{growing} population \cite{SCHUMACHER2009399} which \hl{results} in \hl{the }additional pressure on the environment on the one hand, and \hl{gradual} increase of the built-up area on the other. The confluence of the Blue Nile and White Nile forming its major artery of the Greater Nile near the Khartoum \hl{results} in \hl{high} concentration of population in this district. \hl{Further, consequences} of the desertification \hl{drive} people into \hl{Khartoum}, thereby \hl{increasing} the urbanisation\hl{ which} increases \hl{the} environmental pressure in the\hl{ region} \cite{Babiker1982,Mohammed}. The \hl{Atbara} River (\hl{Red} Nile) \hl{has} high mining activities related to \hl{resource exploration }\cite{SASMAZ2020106539}. 

\hl{The} desertification leads to many negative environmental consequences such as land degradation in arid and semi-arid regions of Sudan \cite{Khiry,ELNASHAR2022152925}. \hl{Social} consequences of desertification in Sudan \hl{include} increased \hl{droughts} and associated famine due to affected food production and soil degradation in the Sahelian zone\hl{ }\cite{Subramaniam}. \hl{The} environmental changes  in the past, \hl{highlighted }links between the human occupation and \hl{responses from environment} \cite{HONEGGER2015141}. For instance, deep ploughing of the surface increased the susceptibility of soil to wind erosion, which \hl{leads} to a severe decline in fertility and \hl{formation} of sand dunes \cite{Abdi}. \hl{Other} examples of the effects of the desertification include reduced biological productivity \cite{KHALIFA2018790,Jahelnabi} and decrease in plant biomass, increase of the extremely dry desert areas. \hl{The} cumulative effect of \hl{these} factors \hl{include }the unpredictable and severe droughts, desiccation or aridification and dryland degradation or desertification in the dryland regions \cite{Darkoh}. 

\hl{Climatic} trends are \hl{reflected} in the hydrology of the Nile \hl{which shows} high levels of \hl{the }variability in river flow records \cite{CONWAY200599} and variations in peak discharge \cite{Babiker}. For instance, the increasing trend in natural water storage variation of the Nile \hl{is} detected \hl{in }northern \hl{Sudan} \cite{ABDELMALIK2019103596}. Such fluctuations in the hydrology of \hl{Nile} necessarily affect the vegetation coverage in the banks and surrounding areas\hl{. }Besides\hl{, }the patterns of the vegetation \hl{in the} Khartoum Province are deeply connected to the regional geologic setting \cite{ELSHEIKH201182}, the geomorphic structure and dominating soil types that control the distribution of the specific plants \cite{Mahmoud}. \hl{The} White Nile \hl{is sensitive} to the rainfall and evaporation at lake and swamps \hl{due }to the specific hydrology. \hl{Therefore, its }responses\hl{ }to occasional climatic events can occur over timescales \cite{Sene}. \hl{The analysis of the} vegetation patterns \hl{and }soil erosion \hl{can be used for detecting vulnerable regions prone to desertification} \cite{Bakhit}. \hl{Nevertheless, }despite \hl{the} efforts on water resources conservation and rational exploitation of resources \cite{MOHAMEDALI2003237,w10111579}, the environmental sustainability in Sudan region remains a challenge.

 
%%%%%%%%%%%%%%%%%%%%%%%%%%%%%%%%%%%%%%%%%%

\section{Materials and Methods}

%----------- subsection ----------->
\subsection{Data}

\hl{The satellite images were selected from the United States Geological Survey (USGS) EarthExplorer repository} (\href{https://earthexplorer.usgs.gov/}{https://earthexplorer.usgs.gov/}). \hl{The choice of the data is explained by the availability of the Landsat OLI/TIRS scenes, sensor spectral band pass, high spatial resolution of 30 m across the multispectral channels (1 to 7), open availability of the cloud-free images (0\% of cloudiness) and regular survey which provides the comparable gap between the target years (2013 and 2022) where the images were collected for analysis of changes in vegetation} (Figure \ref{fig03}).

\begin{figure}[H]
\makebox[1.0\linewidth][r]{%
	\begin{subfigure}[b]{.40\textwidth}
		\centering
			\includegraphics[width=0.90\textwidth]{Fig_03a.jpg}
		\caption{2013}
	\end{subfigure}%
	\begin{subfigure}[b]{.40\textwidth}
		\centering
			\includegraphics[width=0.90\textwidth]{Fig_03b.jpg}
		\caption{2018}
	\end{subfigure}%
	\begin{subfigure}[b]{.40\textwidth}
		\centering
			\includegraphics[width=0.90\textwidth]{Fig_03c.jpg}
		\caption{2022}
	\end{subfigure}%
}
\caption{Landsat 8 OLI/TIRS image of White and Blue Nile in Khartoum, Sudan, in natural RGB colours in December: (\textbf{a}) 20.12.2013, (\textbf{b}) 18.12.2018, (\textbf{c}) 21.12.2022}\label{fig03}
\end{figure}

\hl{Among the available Landsat produces, the Landsat OLI/TIRS sensor was selected, as it has better technical and spectral characteristics compared to older sensors. Thus, compared to the Landsat-7 ETM+, the Landsat 8-9 OLI/TIRS has narrower spectral bands, higher calibration, finer radiometric resolution and geometry, better signal to noise parameters} \cite{ROY201657}. \hl{This results in a more narrow and precise bandwidth in red and NIR channels which differs it from the older sensors. For instance, the Landsat 8-9 OLI/TIRS sensor has 0.64-0.67 \textmu m in the Red band (Band 4), and 0.85-0.88 \textmu m in the NIR (Band 5), while the Landsat 7 ETM+ sensor has coarser range of 0.63-0.69 \textmu m in the Red (Band 3) and 0.77-0.90 \textmu m in NIR (Band 4). The spectral characteristics of the older Landsat products (Landsat 4-5 TM and Landsat 1-5 MSS) has even coarser parameters. Since the Landsat 8 OLI/TIRS was launched on February 11, 2013, there is no earlier images, and the earliest possible date for the cloud-free image on target month of December was 20/12/2013. The images were taken within the 5 year gap between 12/2013 and 12/2018, and 4 year gap between 12/2018 and 12/2022 which gives the comparable intervals of the periods between the scenes.}

The summary of the data is presented in Table \ref{tab1}. All \hl{the }images have 0\% cloudiness \hl{and Path/Row ID 173/49}. \hl{The Sensor ID is common for all the scenes: Landsat 8-9 OLI/TIRS (Operational Land Imager and Thermal Infrared Sensor), Collection 2 Level-2. Image source: the USGS} \cite{EROS}. \hl{The major technical and geodetic metadata common for all the images are the following: Datum and Ellipsoid WGS84; Product Map Projection L1 -- Universal Transverse Mercator (UTM); UTM Zone 36; Sensor ID OLI TIRS; Landsat Processing Software Version LPGS\_16.2.0. The full list of metadata is provided in Table} \ref{tabApp}.

\begin{table}[H]
\scriptsize
\caption{Satellite images used for computing the VIs: Landsat-8 OLI/TIRS collected from the USGS \textsuperscript{1}.\label{tab1}}
%	\begin{adjustwidth}{-\extralength}{0.5cm}
%		\newcolumntype{C}{>{\centering\arraybackslash}X}
		\begin{tabularx}{\linewidth}{ p{1.6cm} | p{1.3cm} | p{6.6cm} | p{2.2cm} }
			\toprule
			\textbf{Date} & \textbf{Spacecraft} & \textbf{Landsat Product ID} & \textbf{Scene ID} \\
			\midrule
			2013/12/20 & Lands. 8 & \makecell{LC08\_L2SP\_173049\_20131220\_20200912\_02\_T1} & LC81730492013354LGN01 \\
			2018/12/18 & Lands. 8 & \makecell{LC08\_L2SP\_173049\_20181218\_20200830\_02\_T1} & LC81730492018352LGN00 \\
			2022/12/21 & Lands. 9 & \makecell{LC09\_L2SP\_173049\_20221221\_20221224\_02\_T1} & LC91730492022355LGN00 \\		
			\bottomrule
		\end{tabularx}
%	\end{adjustwidth}
%	\noindent{\footnotesize{\textsuperscript{1} }}
\end{table}

%----------- subsection ----------->
\subsection{Methodology}

\hl{With this study, we process the satellite images using R version 4.2.2 (R Core Team, 2020) and its packages for the computation of VIs and spatial data analysis}\cite{RCoreTeam}. The extraction of \hl{the }VIs for analyses of vegetation health by traditional software, which is conceptually straightforward due to raster algebra, may lead nevertheless to a consequent computing time. \hl{The VIs are used as robust and reliable indicators of plant health and greenness because they indicate at the level of leaf chlorophyll in plants. In such a way, the VIs reflect the functional features of canopy and parameters of vegetation coverage essential for environmental and climate analytics. }

\hl{Over the course of the past 20 years, there has been a significant progress in programming applications in environmental studies towards automation methods in image processing, such as R scripts. Scripting methods in geoinformatics have undergone a strong evolution during the last decades and now they generally result in a higher accuracy, diverse computational approaches in a variety of libraries and pre-programmed packages for data processing. Such technical performance is beneficial for image processing and spatial data analysis compared to traditional GIS methods }\cite{6049966,Lemenkova202301,6910592,Lemenkova202228,9323708,1027236}. Considering \hl{such} advantages of the programming languages in automation of geospatial data analysis, the algorithms of R present a faster and more effective approach to implement by scripts for image processing \hl{as we demonstrated in this paper}. 

\hl{In this study, we compute the VIs and analyse the maps based on these indices for the years between 2013 and 2022. }In a first step, the necessary libraries 'terra' \hl{version 1.7-18}, 'RColorBrewer' \hl{version 1.1-3} \cite{Neuwirth}, 'Hmisc' \hl{version 5.0-1} \cite{Hmisc} and 'pals' \hl{version 1.7 }\cite{Wright} were activated and the working folder is defined. The 'terra' library is the key package which was used for raster algebra and calculation of the VIs, while the other libraries are used as auxiliary packages for graphical refinements and proper visualization of the plots. Then, we defined the function of the VI, different for each of the VIs using relevant formulae. Finally, we uploaded the necessary Landsat bands in the active folder and performed calculation of each of the VIs by R scripts, as presented below for each corresponding index.

\hl{To compute the VIs, we used the 'terra' R package approach which allows assessment of the vegetation patterns using computational formulae to study vegetation greenness and to assess the increase of desert areas. There are two important advantages of the R approach with regards to estimation of vegetation health using histograms showing data distribution for analysis of the phenological changes in canopies: i) the embedded map algebra is adjusted to describe any type of VIs using function ($vi <- function(img, k, i)$) (see the R scripts in Appendix) that enables to vary Landsat bands for computational purposes in different formulae; and ii) it calculates the data frequency distribution showing the number of pixels that correspond to higher greenness, from which the trends to desertification can be assessed using image analysis for various years.} \hl{The full scripts used for plotting the Figures }\ref{fig01}, \ref{fig04}, \ref{fig05}, \ref{fig06}, \ref{fig07} and \ref{fig08} \hl{are provided in the GitHub repository:} \href{https://github.com/paulinelemenkova/Mapping_Sudan_Scripts}{https://github.com/paulinelemenkova/Mapping\_Sudan\_Scripts}

%----------- subsection ----------->
\subsection{Calculating Vegetation Indices}

\subsubsection{Normalized Difference Vegetation Index (NDVI)}
The Normalized Difference Vegetation Index (NDVI) \cite{Rouse} is calculated as a \hl{ratio} between the sum and the difference of the two Landsat bands - \hl{Near-Infrared} (NIR) and Red. \hl{The NDVI is the most well know and widely used ecological indicator for assessing vegetation conditions due to the optimally created formula that includes the spectral properties of Red and NIR bands. }\hl{Thus, compared to the other bands, the} content of chlorophyll in healthy vegetation is reflected by the \hl{NIR} band much more\hl{. }At the same time, it absorbs Red light. Therefore, using these two bands in a well-established combination \hl{that }gives robust results indicating the presence and amount of chlorophyll in leaves. Such effectiveness of the NDVI resulted in its numerous applications for the environmental analysis \cite{w13192707,rs12244119,GHEBREZGABHER2020249,RULINDA2012169}. \hl{The} values of the NDVI \hl{vary within the range} from -1 to +1 \hl{where} negative values \hl{indicate} water, lower values above 0 show bare soil and land, \hl{and }mediocre values (around 1.30-0.50) signify sparsely vegetated areas\hl{. Higher }values \hl{around} +1 identify dense vegetation with healthy canopy and green foliage. The NDVI is computed using Equation \ref{eq:1}:

\begin{equation}\label{eq:1}
NDVI=\dfrac{(NIR-Red)}{(NIR+Red)}
\end{equation}

Using R\hl{, }we computed the NDVI using script \hl{shown} in \hl{the Appendix, }Listing \ref{lst01}. 

\subsubsection{Green Normalized Difference Vegetation Index (GNDVI)}

The GNDVI index is an updated and modified NDVI to remotely sense the presence and vitality of vegetation. Compared to \hl{the }NDVI, it also uses \hl{the }NIR Band for computation but replaces the \hl{Red} Band (i.e., Band 4 for the Landsat 8-9 OLI/TIRS with diapason of 640-670 nm of wavelength) \hl{by the }Green Band (Band 3 for the Landsat 8-9 OLI/TIRS with 530-590 nm wavelength) from the visible spectrum\hl{.} Therefore, the GNDVI estimates the chlorophyll content more precisely compared to the NDVI and enables to analyse photo synthetic activity in plants. The applications of the GNDVI are mostly intended at detecting wilted, ill or aging plants and monitoring plant stress. 

\hl{Moreover, the GNDVI} is suitable to estimate \hl{the }nitrogen content in \hl{the }plant leaves, as well as\hl{ }monitor mature vegetation coverage in the final stages of the crop cycle, or vegetation with dense canopy. Similar to the NDVI, the range of values in the GNDVI also varies between -1 and 1, where negative values are associated with the areas of water or bare soil, while higher values indicate healthy vegetation. \hl{In this study, the GNDVI was computed due to its technical and practical advantages. Thus, its advantages include the low number of required spectral bands (only NIR and Green, as indicated in the formula of Equation} \ref{eq:5}) \hl{which results in easy computation in R, and reliability in the analysis of vegetation health due to the more precise estimation of the chlorophyll content compared to the NDVI. }The formula for the GNDVI \hl{used for computation }is given in Equation \ref{eq:5}:

\begin{equation}\label{eq:5}
	GNDVI=\dfrac{(NIR-Green)}{(NIR+Green)}
\end{equation}

For the Landsat 8-9 OLI/TIRS bands, the Green channel corresponds to the Band 3 and Red channel -- to the Band 4. The adjustment so the GNDVI are based on the maximum sensitivity of reflectance of leaves since the reflectance near 670 nm of the Red Band is not sensitive to chlorophyll concentration due to the saturation of the relationship between absorptions and chlorophyll concentration \cite{GITELSON1996289}. Therefore, the replacement of the Red Band by the Green Band results in a more accurate to estimate the concentration of chlorophyll and the rate of photosynthesis. Using R syntax, we computed the GNDVI \hl{using script presented in the Appendix, } Listing \ref{lst02}. 

\subsubsection{Normalized Difference Water Index (NDWI)}

Developed by \hl{Gao in 1996 }\cite{GAO1996257} as a complementary to the NDVI, the NDWI is strongly related to the plant water content as it is sensitive to liquid water content of leaves \hl{based on} the ratio between the difference and sum of NIR and SWIR. The second version of the NDWI developed by \cite{rs5073544} uses a Green band instead of SWIR which can be used to detect surface open water areas as well as evaluate the level of turbidity. In this case, the NDWI is computed using the ratio between the difference and sum of the NIR and visible Green Bands \cite{ZHENG2021129488}. \hl{In our study we used the first version developed by Gao following }the Equation \ref{eq:2}:

\begin{equation}\label{eq:2}
	NDWI=\dfrac{(NIR-SWIR)}{(NIR+SWIR)}
\end{equation}

The advantages of the NDWI is that it lessens the effects from the reflectance of soil and vegetation foliage thus is suitable for sandy areas such as in Sahelian Sudan. A special value of the NDWI for environmental monitoring is that it can be used for indicating areas with water deficit \cite{TENG2021107260}. \hl{Therefore, the NDWI presents a useful and precise tool in the assessment of the health of vegetation in arid areas prone to droughts, such as Sudanian Sahel, through computing the moisture level in plants as detected water content in leaves. }Thus, during \hl{the }drought periods, foliage is strongly affected by \hl{the }water deficit which results in plant illness. Specifically for \hl{the }agricultural areas of Khartoum\hl{, }it leads to the crop failure and decreases \hl{the }production \hl{which affects the} harvest. Detecting \hl{the }plant water stress is essential to evaluate \hl{the }current state of plants and make a prognosis for future development, \hl{e.g.,} by indicating \hl{the }areas that \hl{require extra} irrigating. Using R library 'terra', we computed the NDWI \hl{using script presented in the Appendix in Listing} \ref{lst03}.

\subsubsection{Optimized Soil Adjusted Vegetation Index (OSAVI)}

Originally developed by \cite{Rondeaux}, the Optimized Soil Adjusted Vegetation Index (OSAVI) \hl{is well adjusted to monitor vegetation health }specifically in the regions with low vegetation coverage, such as Sahelian region of Sudan \hl{due to corrections of bare soil and desert areas}. \hl{Specifically, the OSAVI }uses a defined value of 0.16 as a factor to adjust the canopy background as this value gives higher soil variation compared to the original SAVI\hl{. }At the same time, \hl{the }OSAVI has a higher sensitivity to vegetation coverage that covers more than a half of the area in percentage. \hl{Due to such characteristics, the OSAVI is suitable to monitor the existing productivity changes in agricultural areas near the Khartoum and along the banks of the Nile River. Moreover, the corrections fro the reflectance from bare soil and lands enables the detect areas prone to droughts using comparison of the multi-temporal satellite images. Therefore, }the application\hl{ }of \hl{the }OSAVI is especially suitable in the regions with sparse vegetation, as in \hl{the }southern Sudan where soil is often visible through the canopy of vegetation. The OSAVI is computed using Equation \ref{eq:4}:

\begin{equation}\label{eq:4}
	OSAVI=\dfrac{(NIR-Red)}{(NIR+Red+0.16)}
\end{equation}

In R language, we computed the OSAVI \hl{by script showed in the Appendix, Listing }\ref{lst04}.

\subsubsection{Infrared Percentage Vegetation Index (IPVI)}
The values of the Infrared Percentage Vegetation Index (IPVI) depend on the chlorophyll content in leaves. Developed originally by \cite{CRIPPEN199071}, it has always positive values ranging from 0 to +1 and indicates the photosynthetic activity of the canopy cover and total of chlorophyll\hl{ in leaves}. Further, as a modified variation\hl{ of the NDVI, }it is especially suitable to indicate yellowish or shed leaves \hl{in the semi-arid areas such as southern Sudan, due }to the adjusted spectral regions in the formula\hl{. Thus, the }chlorophyll content directly depends on \hl{the }nitrogen level in plants which contributes to \hl{the }greenness of foliage and is reflected in \hl{the }optical characteristics of leaves. \hl{In such way, the }IPVI is suitable for the periods of active vegetation development due to the indication of leaf pigment content and variations showing always positive values above zero, in contrast with the NDVI. \hl{Therefore, the IPVI is useful to indicate the yield characterization and stress level in leaves in the regions affected by droughts in the semi-arid ares to the norther of Khartoum. In this way, the IPVI is a robust indicator of the drought severity in southern Sudan, and its effects on the arid and semi-arid ecosystems of Sudanian Sahel. Using the comparison of images covering target test area on the given time periods in 2013, 2018 and 2022, the difference in the vegetation patterns shows the variations in health conditions of leaves. }The IPVI is computed using the Equation \ref{eq:3}:

\begin{equation}\label{eq:3}
	IPVI=\dfrac{NIR}{NIR + Red}
\end{equation}

The computation of IPVI using R library 'terra' was done by the code \hl{presented in the Appendix, Listing }\ref{lst05}. \hl{The} application of RStudio \hl{version 2023.03.0+386 } \cite{RStudio} for image processing \hl{included image processing, computing and} plotting \hl{the vegetation indices} index using Landsat 8 OLI/TIRS image. 

%%%%%%%%%%%%%%%%%%%%%%%%%%%%%%%%%%%%%%%%%%
%\newpage
\section{Results}

We provide in Table \ref{tab02} for each image triple their VI values for years 2013, 2018 and 2022, derived from the R report on data analysis, as it appears in the output console. These values were collected and summarised in Table \ref{tab02} for comparison of variations in the VI's values. \hl{In }the second to fourth columns we indicate the years of observations and output results obtained by R 'terra' library algorithm. 

\begin{table}[H] 
\caption{Evaluation results of the VI values derived from the Landsat 8 OLI/TIRS images over the Khartoum region area (2013, 2018 and 2022).\label{tab02}}
\newcolumntype{C}{>{\centering\arraybackslash}X}
\begin{tabularx}{\textwidth}{CCCC}
\toprule
\textbf{VI}	& \textbf{2013}	& \textbf{2018} & \textbf{2022}\\
\midrule
	$NDVI_{min}$ & -0.2742168 & -0.2817099 & -0.2791084 \\
	$NDVI_{max}$ & 0.7246868 & 0.5796811 & 0.6044010 \\\hline
	$GNDVI_{min}$ & -0.2387020 & -0.2830099 & -0.2253696 \\
	$GNDVI_{max}$ & 0.6371941 & 0.5690726 & 0.7235179 \\\hline
	$NDWI_{min}$ & -0.4046338 & -0.5656914 & -0.5743572 \\
	$NDWI_{max}$ & 0.3880472 & 0.3764042 & 0.4152763\\\hline
	$OSAVI_{min}$ & -0.2742153 & -0.2817084 & -0.2791067 \\
	$OSAVI_{max}$ & 0.7246833 & 0.5796785 & 0.6043976 \\\hline
	$IPVI_{min}$ & 0.1376566 & 0.2101595 & 0.1977995 \\
	$IPVI_{max}$ & 0.6371084 & 0.6408550 & 0.6395542 \\
\bottomrule
\end{tabularx}
\end{table}

\hl{For each VI, the minimal and maximal values were computed using Landsat images taken on three dates: 20.12.2013, 18.12.2018 and 21.12.2022 to assess the range of data. These data show the variations indicating changes in vegetation health. The maximal values indicate the healthy green vegetation which corresponds to greenness due to the chlorophyll content in leaves. The changes between the 2013 and 2018 show the decline for the values in the following indices: NDVI, GNDVI, NDWI, OSAVI and a vey slight increase in the IPVI. The analysis of data on 2018 and 2022 shows that there is a slight rise in values for the same indices expect for the IPVI which demonstrates a slight decrease.}

\subsection{Normalized Difference Vegetation Index (NDVI)}
The images of the NDVI (Figure \ref{fig04}) allow the differentiation of vegetation areas \hl{or cultivated areas from} soil in the region southwards from the Khartoum and along the flow of the\hl{ }Greater Nile. \hl{The }negative NDVI values were classified as non-vegetation (i.e. soil and bare land), and positive NDVI values -- as vegetation. The highest NDVI values (0.80 to 1.0) identify forest areas. For the desert areas with dominating sands such as Sudan, the NDVI values do not exceed 0.50. The minimal values determined according to the histograms obtained from the NDVI images for years 2013, 2018 and 2022 were 0.27, 0.28 and 0.28 (Table \ref{tab02}), whereas, the highest determined values were 0.72, 0.58 and 0.60 for the same years. \hl{The minimum and maximum values were derived from the computed vegetation indices as a statistical summary of data analysis. They indicate the range in greenness in the vegetation: the higher values correspond to the healthy green vegetation while lower values show bare land. The decrease in higher values shows the decline in green and healthy plants}. 

\begin{figure}[H]
\makebox[0.99\linewidth][r]{%
	\begin{subfigure}[b]{.55\textwidth}
		\centering
			\hspace{30pt}\includegraphics[width=0.95\textwidth]{Fig_04a.jpg}
		\caption{}
	\end{subfigure}\hspace*{\fill}%
	\begin{subfigure}[b]{.55\textwidth}
		\centering
			\hspace{30pt}\includegraphics[width=0.95\textwidth]{Fig_04b}
		\caption{}
	\end{subfigure}%
}\\
\vfill \vspace{1mm}
\makebox[0.99\linewidth][r]{%
	\begin{subfigure}[b]{.55\textwidth}
		\centering
			\hspace{30pt}\includegraphics[width=0.95\textwidth]{Fig_04c.jpg}
		\caption{}
	\end{subfigure}%
	\begin{subfigure}[b]{.55\textwidth}
		\centering
			\hspace{30pt}\includegraphics[width=0.95\textwidth]{Fig_04d}
		\caption{}
	\end{subfigure}%
}\\
\vfill \vspace{1mm}
\makebox[0.99\linewidth][r]{%
	\begin{subfigure}[b]{.55\textwidth}
		\centering
			\hspace{30pt}\includegraphics[width=0.95\textwidth]{Fig_04e.jpg}
		\caption{}
	\end{subfigure}%
	\begin{subfigure}[b]{.55\textwidth}
		\centering
			\hspace{30pt}\includegraphics[width=0.95\textwidth]{Fig_04f}
		\caption{}
	\end{subfigure}%
}
\caption{NDVI computed from the Landsat 8 OLI/TIRS images on confluence of the White Nile and Blue Nile, Sudan for three years (always December): (\textbf{a}) 2013, (\textbf{b}) 2018, (\textbf{c}) 2022, (\textbf{d}) histogram for 2013, (\textbf{e}) histogram for 2018, (\textbf{f}) histogram for 2022.}\label{fig04}
\end{figure}

The pixels that were not classified to either the soil, water or vegetation classes, were identified as sand desert areas (mostly beige-coloured in the images in Figure \ref{fig04}) pixels due to low NDVI values around 0.2 or less. The lowest NDVI values less then 0.1 are classified as rocky surface and bare land. Light ivory colour corresponds to the patterns of the dry creeks or beds of wadi ephemeral streams. The minimal and maximal values of the NDVI varied for years 2013, 2018 and 2022 as follows: from -0.2742168 to 0.7246868 in 2013; from -0.2817099 to 0.5796811 in 2018, and from -0.2791084 to 0.6044010 in 2022. As one can note, the maximal values decreased from 2013, which indicated the general trend in desertification, however, fro, 2018  to 2022, there is a slight increase again in the highest values which can be explained by better climate setting and fluctuations in Nile hydrology in 2022 compared to 2018.

\subsection{Green Normalized Difference Vegetation Index (GNDVI)}

The highest values of the GNDVI did not exceed 0.72 for the area of Sudan due to the scarce vegetation coverage and dominating desert areas, Figure \ref{fig05}. 

\begin{figure}[H]
\makebox[0.99\linewidth][r]{%
	\begin{subfigure}[b]{.55\textwidth}
		\centering
			\includegraphics[width=0.95\textwidth]{Fig_05a}
		\caption{2013}
	\end{subfigure}%
	\begin{subfigure}[b]{.55\textwidth}
		\centering
			\includegraphics[width=0.95\textwidth]{Fig_05b}
		\caption{2013: histogram}
	\end{subfigure}%
}\\
\vfill \vspace{1mm}
\makebox[0.99\linewidth][r]{%
	\begin{subfigure}[b]{.55\textwidth}
		\centering
			\includegraphics[width=0.95\textwidth]{Fig_05c}
		\caption{2018}
	\end{subfigure}%
	\begin{subfigure}[b]{.55\textwidth}
		\centering
			\includegraphics[width=0.95\textwidth]{Fig_05d}
		\caption{2018: histogram}
	\end{subfigure}%
}\\
\vfill \vspace{1mm}
\makebox[0.99\linewidth][r]{%
	\begin{subfigure}[b]{.55\textwidth}
		\centering
			\includegraphics[width=0.95\textwidth]{Fig_05e}
		\caption{2022}
	\end{subfigure}%
	\begin{subfigure}[b]{.55\textwidth}
		\centering
			\includegraphics[width=0.95\textwidth]{Fig_05f}
		\caption{2022: histogram}
	\end{subfigure}%
}
\caption{GNDVI computed from the Landsat 8 OLI/TIRS images on confluence of the White Nile and Blue Nile, Sudan for three years (always December): (\textbf{a}) 2013, (\textbf{b}) 2018, (\textbf{c}) 2022, (\textbf{d}) histogram for 2013, (\textbf{e}) histogram for 2018, (\textbf{f}) histogram for 2022.}\label{fig05}
\end{figure}

The negative values are associated with the areas of water or bare soil, while higher values indicate healthy vegetation and selected agricultural croplands along the banks of the White, Blue and Greater Niles. \hl{As seen from Figure} \ref{fig05}, the GNDVI captured plant greenness better than the NDVI: thus, changes in plant colour from bright green to yellow on sequence of images for 2013, 2018 and 2022 are depicted in a more refined way in the GNDVI image compared to NDVI. In this case, the variations in shadows of colours indicate the differences in plant vigour and health which in turn point at the irrigated lands with higher level of moisture that can be distinguished against drylands and desert areas. 

Further, comparing the values in the Table \ref{tab02}, the \hl{range} of values for the GNDVI is narrower than that of the NDVI: with minimal -0.24, -0.28 and -0.23 for years 2013, 2018 and 2022, while the highest values are 0.64, 0.57 and 0.72 for the same years. Accordingly, greener plants visualised in the GNDVI, Figure \ref{fig05} have lower values yet are better distinguished in Figure \ref{fig05}. Moreover, the histograms present more diversified range of values for various land classes: water, bare soil and sand, rocky areas of deserts and vegetation with agricultural areas and natural aquatic plants in the Nile rivers. Finally, the slight increase in the GNDVI values indicating photo synthetic activity in leaves is affected by changes in water and N absorption in plant that includes aquatic vegetation in the bank areas of the White Nile. 

\subsection{Normalized Difference Water Index (NDWI)}

The interpretation of the values on the resulting images based on the NDWI (Figure \ref{fig06}) are as follows: the negative values (from -1 up to 0) signify areas with no vegetation or water content, while those above 0 show surfaces with water content where the higher moisture level correlates with higher values of the NDWI, respectively. The NDWI values vary between -1 to +1, which depends on the total amount of water in leaves, type of vegetation and percentage of the foliage coverage. Thus, sparse vegetation in semi-arid regions of the Sudanian Sahel will naturally have a lower value of the NDWI compared to the regions of dominating forests with dense canopy. A closer look at the NDWI images reveals lines contouring patterns of the streams of the dry creeks in the desert area with typical  braided stream patterns and beds of drainage courses, indicating differences in moisture content with most area (Figure \ref{fig06}). 

Such patterns are indicative of moist wadi sediments and associated vegetation strips. The range of values in the NDWI shows slow dynamics in the lowest values that is gradually increasing from -0.40 in 2013, to -0.56 in 2018 and 0.57 in 2022, along with the raising highest values: 0.38 in 2013, 0.38 in 2018 and 0.42 in 2022. The fluctuations in the NDWI values are affected by the increase in water discharge in White Nile River over nine years. Since the NDWI is specifically designed to monitor vegetation in drought affected areas such as Sudan, such variations indicate at the increase in both the content of vegetation water because of the strong chlorophyll absorption in SWIR and the structure in plant canopy. 

\hl{On one hand}, partial vegetation coverage in the desert areas of central Sudan results in soil effects on the NDWI values, especially for the landscapes with interspersed soils and vegetation patterns. Thus, when the leaf layer in the canopy decreases, the absolute values of reflectances in the spectral region of the NDWI are decreased which results in lower values. In general, the range of the NDVI values over the green vegetation areas in Khartoum is narrow, which is reflected in the elongated histogram where most of the values are from -0.20 to +0.20. More specifically, the interpretation of the values of the NDWI follows the ranges: high positive values from +0,2 to +1.0 -- water area; 0.0 to +0,2 -- flooded areas and high humidity, -0,3 to 0.0: moderate drought soil and non-water surfaces, -1.0 to -0.3: drylands, drought and non-water surfaces. 

Our results show that the NDWI performs better on identification of vegetation in the sandy areas  compared to the NDVI and GNDVI, which can also be noted by examination and comparison of on the histograms which have a more gentle curve with detailed range of data values for the case of NDWI for the Sahelian Sudan.

\begin{figure}[H]
\makebox[0.99\linewidth][r]{%
	\begin{subfigure}[b]{.5\textwidth}
		\centering
			\includegraphics[width=0.95\textwidth]{Fig_06a}
		\caption{2013}
	\end{subfigure}%
	\begin{subfigure}[b]{.5\textwidth}
		\centering
			\includegraphics[width=0.95\textwidth]{Fig_06b}
		\caption{2013: histogram}
	\end{subfigure}%
}\\
\vfill \vspace{1mm}
\makebox[0.99\linewidth][r]{%
	\begin{subfigure}[b]{.5\textwidth}
		\centering
			\includegraphics[width=0.95\textwidth]{Fig_06c}
		\caption{2018}
	\end{subfigure}%
	\begin{subfigure}[b]{.5\textwidth}
		\centering
			\includegraphics[width=0.95\textwidth]{Fig_06d}
		\caption{2018: histogram}
	\end{subfigure}%
}\\
\vfill \vspace{1mm}
\makebox[0.99\linewidth][r]{%
	\begin{subfigure}[b]{.5\textwidth}
		\centering
			\includegraphics[width=0.95\textwidth]{Fig_06e}
		\caption{2022}
	\end{subfigure}%
	\begin{subfigure}[b]{.5\textwidth}
		\centering
			\includegraphics[width=0.95\textwidth]{Fig_06f}
		\caption{2022: histogram}
	\end{subfigure}%
}
\caption{NDWI computed from the Landsat 8 OLI/TIRS images on confluence of the White Nile and Blue Nile, Sudan for three years (always December): (\textbf{a}) 2013, (\textbf{b}) 2018, (\textbf{c}) 2022, (\textbf{d}) histogram for 2013, (\textbf{e}) histogram for 2018, (\textbf{f}) histogram for 2022.}\label{fig06}
\end{figure}

%--------------------->
\subsection{Optimized Soil Adjusted Vegetation Index (OSAVI)}

The OSAVI was the highest for 2013 with the maximal reached value of 0.72 (Figure \ref{fig07}), after which it dropped to 0.58 in 2018 indicating the decease in the leaf chlorophyll content reflected in the OSAVI values; since 2018 to 2022 the values stabilised and slightly increased to 0.60 due to the effects from the White Nile hydrology and climatic variations. Being a modified index derived from the NDVI, the OSAVI also has a range of values varying between the -0.28 to +0.72, with negative values indicating water areas of the White and Blue Niles, lower areas slightly above 0 and lower than 0.2 indicating bare land, sandy areas, soil with no or very sparse vegetation, while higher values close to +0.50 showing healthy vegetation and higher chlorophyll content in leaves. 

\begin{figure}[H]
\makebox[0.99\linewidth][r]{%
	\begin{subfigure}[b]{.5\textwidth}
		\centering
			\includegraphics[width=0.95\textwidth]{Fig_07a}
		\caption{2013}
	\end{subfigure}%
	\begin{subfigure}[b]{.5\textwidth}
		\centering
			\includegraphics[width=0.95\textwidth]{Fig_07b}
		\caption{2013: histogram}
	\end{subfigure}%
}\\
\vfill \vspace{1mm}
\makebox[0.99\linewidth][r]{%
	\begin{subfigure}[b]{.5\textwidth}
		\centering
			\includegraphics[width=0.95\textwidth]{Fig_07c}
		\caption{2018}
	\end{subfigure}%
	\begin{subfigure}[b]{.5\textwidth}
		\centering
			\includegraphics[width=0.95\textwidth]{Fig_07d}
		\caption{2018: histogram}
	\end{subfigure}%
}\\
\vfill \vspace{1mm}
\makebox[0.99\linewidth][r]{%
	\begin{subfigure}[b]{.5\textwidth}
		\centering
			\includegraphics[width=0.95\textwidth]{Fig_07e}
		\caption{2022}
	\end{subfigure}%
	\begin{subfigure}[b]{.5\textwidth}
		\centering
			\includegraphics[width=0.95\textwidth]{Fig_07f}
		\caption{2022: histogram}
	\end{subfigure}%
}
\caption{OSAVI computed from the Landsat 8 OLI/TIRS images on confluence of the White Nile and Blue Nile, Sudan for three years (always December): (\textbf{a}) 2013, (\textbf{b}) 2018, (\textbf{c}) 2022, (\textbf{d}) histogram for 2013, (\textbf{e}) histogram for 2018, (\textbf{f}) histogram for 2022.}\label{fig07}
\end{figure}

The OSAVI image shows that the land classes with values between 0.30 and 0.40 representing sparse vegetation, 0.40 to 0.50 representing medium vegetated area typically for agriculture lands and those above 0.50 representing occasional regions with high vegetated area around the Blue Nile. Close examination of the OSAVI and comparison with other indices reveals that the OSAVI classifies pixels on the image more was detailed compared to the NDVI and GNDVI. Thus, it identifies vegetation better using spectral reflectance within each cluster of data range which results in a more accurate detection of the biomass and leaf foliage areas in the regions covered by sandy soils and semi-arid lands such as Sudan. 
%--------------------->
\subsection{Infrared Percentage Vegetation Index (IPVI)}

The majority of values of the Infrared Percentage Vegetation Index (IPVI) for the area of Sudan, Figure \ref{fig08}, lies within a very narrow interval of 0.25 to 0.48 which well corresponds with the dominance of sandy dune areas and bare soils in the surrounding deserts. Note that all the values of IPVI are never negative and always positive due to the formula used in computation, which distinct this index from the previous ones.

\begin{figure}[H]
\makebox[0.99\linewidth][r]{%
	\begin{subfigure}[b]{.55\textwidth}
		\centering
			\includegraphics[width=0.95\textwidth]{Fig_08a}
		\caption{2013}
	\end{subfigure}%
	\begin{subfigure}[b]{.55\textwidth}
		\centering
			\includegraphics[width=0.95\textwidth]{Fig_08b}
		\caption{2013: histogram}
	\end{subfigure}%
}\\
\vfill \vspace{1mm}
\makebox[0.99\linewidth][r]{%
	\begin{subfigure}[b]{.55\textwidth}
		\centering
			\includegraphics[width=0.95\textwidth]{Fig_08c}
		\caption{2018}
	\end{subfigure}%
	\begin{subfigure}[b]{.55\textwidth}
		\centering
			\includegraphics[width=0.95\textwidth]{Fig_08d}
		\caption{2018: histogram}
	\end{subfigure}%
}\\
\vfill \vspace{1mm}
\makebox[0.99\linewidth][r]{%
	\begin{subfigure}[b]{.55\textwidth}
		\centering
			\includegraphics[width=0.95\textwidth]{Fig_08e}
		\caption{2022}
	\end{subfigure}%
	\begin{subfigure}[b]{.55\textwidth}
		\centering
			\includegraphics[width=0.95\textwidth]{Fig_08f}
		\caption{2022: histogram}
	\end{subfigure}%
}
\caption{IPVI computed from the Landsat 8 OLI/TIRS images on confluence of the White Nile and Blue Nile, Sudan for three years (always December): (\textbf{a}) 2013, (\textbf{b}) 2018, (\textbf{c}) 2022, (\textbf{d}) histogram for 2013, (\textbf{e}) histogram for 2018, (\textbf{f}) histogram for 2022.}\label{fig08}
\end{figure}

%%%%%%%%%%%%%%%%%%%%%%%%%%%%%%%%%%%%%%%%%%
\section{Discussion}

\hl{The method of R scripts is advantageous} for repetitive workflow of image processing when several scenes need to be processed sequentially for several years. \hl{Specifically, scripts are highly re-useable which enables to save time through automation of the repetitive tasks. This is especially useful when one needs a computation of several Landsat images to calculate several VIs using formulae applied to the Landsat bands. Further, R scripts eliminate errors when processing the satellite images due to high level of automation. Applying scripts slightly modified for each scene ensures no errors or misclassification of pixels when processing the images which are possible otherwise when working manually in GIS. Finally, scripts facilitate a package-based approach of R that are used as target tools to image processing and other tasks. For instance, scripts call embedded functions to perform parts of computations and data processing. Moreover, scripts provide quick plotting of maps and graphics that visualise the results. This help to highlight the changes in vegetation patterns when maps visualise the VIs on several years for comparison. }

The environment challenges facing the landscapes of Sudan with  semi-arid climate affect the ecosystems. The functionality and structure of the vegetation in Sudan strongly depend upon the hydrology of Nile, especially its southern region near the confluence of the White Nile and the Blue Nile. Moreover, it is controlled by the regional climate setting and variations in rainfall, evaporation, variations in of cloudiness and atmospheric radiation and other climate-related processes. Recent droughts and decrease in rainfall triggered environmental degradation and affected vegetation growth and health due to the lack of humidity in the air. As a consequence, this initiated gradual processes of desertification in Saharan-Sahelian Africa, expansion of drylands and deterioration of the landscapes. 

Nevertheless, there are vast opportunities for sustainable environmental monitoring of Sudan through integration of the remotely sensed data and advanced programming technologies. In this paper, we have proposed the application of R language for computation of several vegetation indices using library 'terra' and satellite images Landsat 8-9 OLI/ITIRS. Compared to the traditional GIS approaches, the proposed method has demonstrated high efficiency and effectiveness for geospatial data processing aimed at environmental analysis. In the performed experiments on calculation of VIs over the Landsat 8 OLI/TIRS data, our study has shown that, in most cases of the indices, a general trend in decrease of VI values is notable for data on 2013, 2018 and 2022, as shown in Table \ref{tab02} and visualised in Figures \ref{fig04}, \ref{fig05}, \ref{fig06}, \ref{fig07} and \ref{fig08}. This is clearly visible from the results obtained using the estimated parameters on indices NDVI, OSAVI, IPVI. 

Using three satellite images Landsat 8-9 OLI/ITIRS collected for the region of Khartoum area, we detected a slight decrease in vegetation greenness from 2013 to 2022 based on the comparison of the dynamics of the satellite-derived VIs. For each image, the computation of five of the presented VIs has been performed for three corresponding years by applied R scripts. The visualization of these indices demonstrated a function of vegetation health, greenness and vigour in the conditions of the Sahelian climate of Sudan with variations over the analysed years. The proposed research technically implemented in R shows promising results on detecting anomalies in vegetation conditions in southern Sudan, as the VIs provide essential information on plant greenness reflected as chlorophyll content. The chlorophyll content related to the water content in leaves is in turn controlled by regional climatic setting and hydrological processes of the Nile River and its tributaries - the White and Blue Niles. 

We have addressed the problem of satellite image processing for computing a series of vegetation indices for the regions prone to desertification in the Sahelian area of Sudan. Our approach represents each Landsat 8-9 OLI/TIRS image with a derived five diverse vegetation indices tuned for various environmental purposes and compared for years 2013, 2018 and 2022: detecting dense vegetation with healthy canopy and green foliage (NDVI), estimation of the canopy background in the regions with sparse vegetation and low percentage of foliage coverage as typical for Sudan (OSAVI), evaluating the maturity of the canopy based on the leaf pigment content and nitrogen level in plants (IPVI), evaluation of concentration of chlorophyll in leaves and the rate of photosynthesis (GNDVI), indicating areas with water deficit (NDWI). In this way, the comparative analysis of indices visualised for various periods can be used to identify vegetation conditions aimed at environmental monitoring of Sahel region of Northeast Africa.

We presented further insights into raster image processing by R scripting and optimization methods for a time series analysis of several scenes. The computational workflow has demonstrated that the proposed algorithms of R perform is far better in terms of the processing and analysis of spatial datasets. The experiments have demonstrated that scientific scripting by R plays an important role in geospatial image processing, visualization and analysis since it enables integrating various components of libraries including graphical packages used for mapping and plotting and packages for spatial data processing which can be integrated into a scripting workflow.  

%%%%%%%%%%%%%%%%%%%%%%%%%%%%%%%%%%%%%%%%%%
\section{Conclusions}

\hl{In this paper, we highlighted that the performance of the R packages is effective for a framework of satellite image processing and enables to compute the vegetation indices. With the presented maps of the VIs made based on the Landsat OLI/TIRS images, we observe the changes in the vegetation patterns in the test area. We further showed that there is a trend to the growth and greenness of vegetation over the period of 2013 to 2022 in the study area in the confluence of the White and Blue Niles in southern Sudan, which is caused by the rising temperature and decreased precipitation due to the climatic and environmental impacts effects of desertification.}

The Landsat 8-9 OLI/TIRS images contain a complex image texture with distinct values of pixels due to different spectral reflectance of vegetation and other land cover types. Five VIs of vegetation distribution analysis have been applied and brought to practical applicability by means of R language. \hl{Specifically, we tested five VIs for the Landsat 8-9 OLI/TIRS images and compared the results for the following indices: 1) Normalized Difference Vegetation Index (NDVI), 2) Normalized Difference Water Index (NDWI), 3) Infrared Percentage Vegetation Index (IPVI), 4) Optimized Soil Adjusted Vegetation Index (OSAVI) and 5) Green Normalized Difference Vegetation Index (GNDVI). }The computational workflow of R libraries included extraction of image sensor DN values of pixels, computing the data range with min/max values, classifying the raster object and visualising the image using selected colour schemes. Within all these methods, the velocity of data processing can be estimated as high compared to conventional GIS approaches, which offers additional possibilities of programming for geospatial data analysis and image processing. In such a way, we demonstrated the machine-based extracting information from satellite images by R programming. We also shows the effectiveness of the remote sensing data as a source of information derived from Earth observation satellites for regional studies focused on Northeastern Africa.

The art of image analysis in Earth sciences has made big progress due to the programming applications and scripts that significantly automate the routine of data processing. In the near future we can expect the possibility of stream processing applied for large collections of satellite images as a series of data for operative monitoring. Thus, further environmental applications of image processing could include, for instance, monitoring burned areas in savannah fires using satellite scenes processed at the time they are generated by sensors. Further, the use of scripts for image processing opens the way to geospatial analysis of diverse objects and structures of very high complexity that constituting the landscapes of the Earth. This includes the detailed diversification of the land cover types, vegetation associations and plant communities that can be identified on images.

In monitoring of the Earth's landscapes, not only the topographic data should be used, but also the environmental characteristics that reflect the complex interactions from climate and hydrological parameters, and vegetation responses to the external climatic effects. These can be derived from the satellite images such as Landsat 8-9 OLI/TIRS using advanced methods of image processing as demonstrated in this paper. Future research will be dedicated to developing further approaches using scripting algorithms aimed at dealing with satellite images to detect other parameters of land surface and vegetation coverage for environmental mapping such as brightness and saturation indices,  evaluating other vegetation indices for detecting the salinity of soil in arid lands of Africa. Such characteristics reflect the climatic-vegetation interactions that can be used for operative environmental monitoring based on the spaceborne data.  

%%%%%%%%%%%%%%%%%%%%%%%%%%%%%%%%%%%%%%%%%%
\authorcontributions{Supervision, conceptualization, methodology, software, resources, funding acquisition, and project administration, O.D.; writing---original draft preparation, methodology, software, data curation, visualization, formal analysis, validation, writing---review and editing, and investigation, P.L. All authors have read and agreed to the published version of the manuscript.}

\funding{The publication was funded by the Editorial Office of Atmosphere, Multidisciplinary Digital Publishing Institute (MDPI), by providing 100\% discount for the APC of this manuscript. This project was supported by the Federal Public Planning Service Science Policy or Belgian Science Policy Office, Federal Science Policy - BELSPO (B2/202/P2/SEISMOSTORM).}

\institutionalreview{Not applicable.}

\informedconsent{Not applicable.}

\dataavailability{Not applicable} 

\acknowledgments{The authors thank the reviewers for reading and review of this manuscript.}

\conflictsofinterest{The authors declare no conflict of interest.} 

\abbreviations{Abbreviations}{
The following abbreviations are used in this manuscript:\\

\noindent 
\begin{tabular}{@{}ll}
GMT & Generic Mapping Tools \\
GNDVI & Green Normalized Difference Vegetation Index \\
IPVI & Infrared Percentage Vegetation Index \\
Landsat 8-9 OLI/TIRS & Landsat 8-9 Operational Land Imager (OLI) and Thermal Infrared (TIRS) \\
NDVI & Normalized Difference Vegetation Index \\
NDWI & Normalized Difference Water Index \\
OSAVI & Optimized Soil Adjusted Vegetation Index \\
SAVI & Soil Adjusted Vegetation Index \\
USGS & United States Geological Survey \\
\hl{UTM} & \hl{Universal Transverse Mercator} \\
VIs & Vegetation Indices \\
\hl{WGS84} & \hl{World Geodetic System 1984}
\end{tabular}
}

%% Optional
\appendixtitles{no} % Leave argument "no" if all appendix headings stay EMPTY (then no dot is printed after "Appendix A"). If the appendix sections contain a heading then change the argument to "yes".
\appendixstart
\appendix
\section[\appendixname~\thesection]{Metadata on the Landsat 8-9 OLI/TIRS images (USGS EarthExplorer).}

\begin{table}[H]
\scriptsize
\caption{Metadata on the Landsat OLI/TIRS satellite images used in this study \textsuperscript{1}.\label{tabApp}}
\begin{adjustwidth}{-\extralength}{0cm}
\begin{tabularx}{\fulllength}{l|l|l|l}
    \hline
        \textbf{Data Set Attribute} & \textbf{Attributes for image on 2022/12/21} & \textbf{Attributes for 2018/12/18} & \textbf{Attributes for 2013/12/20} \\ \hline
        Landsat Scene Identifier & LC91730492022355LGN01 & LC81730492018352LGN00 & LC81730492013354LGN01 \\ \hline
        Date Acquired & 2022/12/21 & 2018/12/18 & 2013/12/20 \\ \hline
        Collection Category & T1 & T1 & T1 \\ \hline
        Collection Number & 2 & 2 & 2 \\ \hline
        WRS Path & 173 & 173 & 173 \\ \hline
        WRS Row & 49 & 49 & 49 \\ \hline
        Target WRS Path & 173 & 173 & 173 \\ \hline
        Target WRS Row & 49 & 49 & 49 \\ \hline
        Nadir/Off Nadir & NADIR & NADIR & NADIR \\ \hline
        Roll Angle & -0.001 & 0.000 & 0.000 \\ \hline
        Date Product Generated L2 & 2023/03/17 & 2020/08/30 & 2020/09/12 \\ \hline
        Date Product Generated L1 & 2023/03/17 & 2020/08/30 & 2020/09/12 \\ \hline
        Start Time & 2022-12-21 08:09:26 & 2018-12-18 08:08:54.736348 & 2013-12-20 08:10:29.988456 \\ \hline
        Stop Time & 2022-12-21 08:09:58 & 2018-12-18 08:09:26.506347 & 2013-12-20 08:11:01.758452 \\ \hline
        Station Identifier & LGN & LGN & LGN \\ \hline
        Day/Night Indicator & DAY & DAY & DAY \\ \hline
        Land Cloud Cover & 0.00 & 0.00 & 0.00 \\ \hline
        Scene Cloud Cover L1 & 0.00 & 0.00 & 0.00 \\ \hline
        Ground Control Points Model & 511 & 521 & 639 \\ \hline
        Ground Control Points Version & 5 & 5 & 5 \\ \hline
        Geometric RMSE Model & 7.263 & 6.387 & 4.378 \\ \hline
        Geometric RMSE Model X & 5.055 & 4.612 & 2.951 \\ \hline
        Geometric RMSE Model Y & 5.216 & 4.418 & 3.234 \\ \hline
        Processing Software Version & LPGS\_16.2.0 & LPGS\_15.3.1c & LPGS\_15.3.1c \\ \hline
        Sun Elevation L0RA & 44.22152569 & 44.38827439 & 44.40055924 \\ \hline
        Sun Azimuth L0RA & 148.68774152 & 148.91482828 & 149.05309430 \\ \hline
        TIRS SSM Model & N/A & FINAL & ACTUAL \\ \hline
        Data Type L2 & OLI\_TIRS\_L2SP & OLI\_TIRS\_L2SP & OLI\_TIRS\_L2SP \\ \hline
        Sensor Identifier & OLI\_TIRS & OLI\_TIRS & OLI\_TIRS \\ \hline
        Satellite & 9 & 8 & 8 \\ \hline
        Product Map Projection L1 & UTM & UTM & UTM \\ \hline
        UTM Zone & 36 & 36 & 36 \\ \hline
        Datum & WGS84 & WGS84 & WGS84 \\ \hline
        Ellipsoid & WGS84 & WGS84 & WGS84 \\ \hline
        Scene Center Lat DMS & "15°54'03.06""N" & "15°54'03.56""N" & "15°54'03.42""N" \\ \hline
        Scene Center Long DMS & "33°06'32.69""E" & "33°07'35.51""E" & "33°05'59.28""E" \\ \hline
        Corner Upper Left Lat DMS & "16°56'34.84""N" & "16°56'44.88""N" & "16°56'44.41""N" \\ \hline
        Corner Upper Left Long DMS & "32°02'24.68""E" & "32°03'25.49""E" & "32°01'44.08""E" \\ \hline
        Corner Upper Right Lat DMS & "16°56'30.70""N" & "16°56'40.13""N" & "16°56'40.63""N" \\ \hline
        Corner Upper Right Long DMS & "34°10'42.96""E" & "34°11'43.87""E" & "34°10'12.61""E" \\ \hline
        Corner Lower Left Lat DMS & "14°50'38.76""N" & "14°50'39.01""N" & "14°50'38.58""N" \\ \hline
        Corner Lower Left Long DMS & "32°03'00.32""E" & "32°04'00.55""E" & "32°02'20.18""E" \\ \hline
        Corner Lower Right Lat DMS & "14°50'35.16""N" & "14°50'34.87""N" & "14°50'35.34""N" \\ \hline
        Corner Lower Right Long DMS & "34°09'59.18""E" & "34°10'59.41""E" & "34°09'29.09""E" \\ \hline
        Scene Center Latitude & 15.90085 & 15.90099 & 15.90095 \\ \hline
        Scene Center Longitude & 33.10908 & 33.12653 & 33.09980 \\ \hline
        Corner Upper Left Latitude & 16.94301 & 16.94580 & 16.94567 \\ \hline
        Corner Upper Left Longitude & 32.04019 & 32.05708 & 32.02891 \\ \hline
        Corner Upper Right Latitude & 16.94186 & 16.94448 & 16.94462 \\ \hline
        Corner Upper Right Longitude & 34.17860 & 34.19552 & 34.17017 \\ \hline
        Corner Lower Left Latitude & 14.84410 & 14.84417 & 14.84405 \\ \hline
        Corner Lower Left Longitude & 32.05009 & 32.06682 & 32.03894 \\ \hline
        Corner Lower Right Latitude & 14.84310 & 14.84302 & 14.84315 \\ \hline
        Corner Lower Right Longitude & 34.16644 & 34.18317 & 34.15808 \\ \hline
%    \end{tabular}
 \end{tabularx}
 \end{adjustwidth}
\end{table}

\section[\appendixname~\thesection]{R scripts used for computing and plotting the vegetation indices}
\subsection[\appendixname~\thesubsection]{R code for computing the NDVI}

\begin{lstlisting}[language=R,caption=R code for computing the NDVI (here: a case of 2022),style=mystyle,label={lst01}]
library(terra)
library(RColorBrewer)
library(Hmisc)
library(pals)
setwd("/Users/polinalemenkova/Documents/R/52_Image_Processing/Sudan_NDVI")
# 1. Normalized Difference Vegetation Index (NDVI) = (NIR - R) / (NIR + R), i.e., NDVI = (Band 5 - Band 4) / (Band 5 + Band 4).
vi <- function(img, k, i) {
  bk <- img[[k]]
  bi <- img[[i]]
  vi <- (bk - bi) / (bk + bi)
  return(vi)
}
# For Landsat NIR = 5, red = 4.
filenames <- paste0('LC09_L2SP_173049_20221221_20221224_02_T1_SR_B', 1:7, ".tif")
filenames
landsat <- rast(filenames)
landsat
ndvi <- vi(landsat, 5, 4)
options(scipen=10000)
colors <- brewer.rdylgn(100)
plot(ndvi, col=colors, font.main = 1, main = "NDVI for Landsat-8 OLI/TIRS C1 image LC09_L2SP_173049_20221221_20221224_02_T1_SR_B: \nWhite Nile and Blue Nile confluence, Sudan (2022)", cex.main=0.9)
minor.tick(nx = 10, ny = 10, tick.ratio = 0.3)
# Plotting the histogram of the NDVI
hist(ndvi, font.main = 1, main = "NDVI values for Landsat-8 OLI/TIRS C1 image \nLC09_L2SP_173049_20221221_20221224_02_T1_SR_B: White Nile and Blue Nile confluence, Sudan (2022)", xlab = "NDVI", ylab= "Frequency",
    col = "darkgoldenrod1", xlim = c(-0.5, 1),  breaks = 30, xaxt = "n")
axis(side=1, at = seq(-0.6, 1, 0.1), labels = seq(-0.6, 1, 0.1))
minor.tick(nx = 10, ny = 10, tick.ratio = 0.3)
\end{lstlisting}

\subsection[\appendixname~\thesubsection]{R code for computing the GNDVI}
\begin{lstlisting}[language=R,caption=R code for computing the GNDVI (here: a case of 2013),style=mystyle,label={lst02}]
setwd("/Users/polinalemenkova/Documents/R/52_Image_Processing/Sudan_GNDVI")
# Green Normalized Difference Vegetation Index (GNDVI) = (NIR - Green) / (NIR + Green)
# Nir = Band 5, Green = Band 3.
vi <- function(img, k, i) {
  bk <- img[[k]]
  bi <- img[[i]]
  vi <- (bk - bi) / (bk + bi)
  return(vi)
}
# For Landsat 8-9 OLI/TIRS, NIR = 5, Green = 3.
filenames <- paste0('LC08_L2SP_173049_20131220_20200912_02_T1_SR_B', 1:7, ".tif")
filenames
landsat <- rast(filenames)
landsat
gndvi <- vi(landsat, 5, 3)
options(scipen=10000)
colors <- brewer.spectral(100)
plot(gndvi, col=colors, font.main = 1, main = "GNDVI for Landsat-8 OLI/TIRS C1 image LC08_L2SP_173049_20131220_20200912_02_T1_SR_B: \nWhite Nile and Blue Nile confluence, Sudan (2013)", cex.main=0.9)
minor.tick(nx = 10, ny = 10, tick.ratio = 0.3)
# Plotting histogram of the GNDVI
hist(gndvi, font.main = 1, main = "GNDVI values for Landsat-8 OLI/TIRS C1 image \nLC08_L2SP_173049_20131220_20200912_02_T1_SR_B: White Nile and Blue Nile confluence, Sudan (2013)", xlab = "GNDVI", ylab= "Frequency", col = "peachpuff", xlim = c(-0.5, 1),  breaks = 50, xaxt = "n")
axis(side=1, at = seq(-0.6, 1, 0.1), labels = seq(-0.6, 1, 0.1))
minor.tick(nx = 10, ny = 10, tick.ratio = 0.3)
\end{lstlisting}

\subsection[\appendixname~\thesubsection]{R code for computing the NDWI}
\begin{lstlisting}[language=R,caption=R code for computing the NDWI (here: a case of 2018),style=mystyle,label={lst03}]
vi <- function(img, k, i) {
  bk <- img[[k]]
  bi <- img[[i]]
  vi <- (bi - bk) / (bi + bk)
  return(vi)
}
# For Landsat 8-9 OLI/TIRS, NIR = 5, SWIR1 = Band 6.
filenames <- paste0('LC08_L2SP_173049_20181218_20200830_02_T1_SR_B', 1:7, ".tif")
filenames
landsat <- rast(filenames)
landsat
ndwi <- vi(landsat, 5, 6)
options(scipen=10000)
colors <- jet(100)
plot(ndwi, col=colors, font.main = 1, main = "NDWI for Landsat-8 OLI/TIRS C1 image LC08_L2SP_173049_20181218_20200830_02_T1_SR_B: \nWhite Nile and Blue Nile confluence, Sudan (2018)", cex.main=0.9)
minor.tick(nx = 10, ny = 10, tick.ratio = 0.3)
# Plotting histogram of the NDWI
hist(ndwi, font.main = 1, main = "NDWI values for Landsat-8 OLI/TIRS C1 image \nLC08_L2SP_173049_20181218_20200830_02_T1_SR_B: White Nile and Blue Nile confluence, Sudan (2018)", xlab = "NDWI", ylab= "Frequency",
    col = "mediumpurple1", xlim = c(-0.5, 1),  breaks = 50, xaxt = "n")
axis(side=1, at = seq(-0.6, 1, 0.1), labels = seq(-0.6, 1, 0.1))
minor.tick(nx = 10, ny = 10, tick.ratio = 0.3)
\end{lstlisting}

\subsection[\appendixname~\thesubsection]{R code for computing the OSAVI}
\begin{lstlisting}[language=R,caption=R code for computing the OSAVI (here: a case of 2022),style=mystyle,label={lst04}]
setwd("/Users/polinalemenkova/Documents/R/52_Image_Processing/Sudan_OSAVI")
# 1. Optimized Soil Adjusted Vegetation Index (OSAVI) = (NIR - R) / (NIR + R + 0.16), i.e., OSAVI = (Band 5 - Band 4) / (Band 5 + Band 4 + 0.16).
vi <- function(img, k, i) {
  bk <- img[[k]]
  bi <- img[[i]]
  vi <- (bk - bi) / (bk + bi + 0.16)
  return(vi)
}
# For Landsat 8-9 OLI/TIRS, NIR = 5, red = 4.
filenames <- paste0('LC09_L2SP_173049_20221221_20221224_02_T1_SR_B', 1:7, ".tif")
filenames
landsat <- rast(filenames)
landsat
osavi <- vi(landsat, 5, 4)
options(scipen=10000)
colors <- kovesi.diverging_rainbow_bgymr_45_85_c67(100)
plot(osavi, col=colors, font.main = 1, main = "OSAVI for Landsat-8 OLI/TIRS C1 image LC09_L2SP_173049_20221221_20221224_02_T1_SR_B: \nWhite Nile and Blue Nile confluence, Sudan (2022)", cex.main=0.9)
minor.tick(nx = 10, ny = 10, tick.ratio = 0.3)
# Plotting the histogram of the OSAVI
hist(osavi, font.main = 1, main = "OSAVI values for Landsat-8 OLI/TIRS C1 image \nLC09_L2SP_173049_20221221_20221224_02_T1_SR_B: White Nile and Blue Nile confluence, Sudan (2022)", xlab = "OSAVI", ylab= "Frequency", col = "dodgerblue", xlim = c(-0.5, 1),  breaks = 30, xaxt = "n")
axis(side=1, at = seq(-0.6, 1, 0.1), labels = seq(-0.6, 1, 0.1))
minor.tick(nx = 10, ny = 10, tick.ratio = 0.3)
\end{lstlisting}

\subsection[\appendixname~\thesubsection]{R code for computing the IPVI}
\begin{lstlisting}[language=R,caption=R code for computing the IPVI (here: a case of 2013),style=mystyle,label={lst05}]
library(terra)
library(RColorBrewer)
library(Hmisc)
library(pals)
#
setwd("/Users/polinalemenkova/Documents/R/52_Image_Processing/Sudan_IPVI")
# Infrared Percentage Vegetation Index (IPVI) = NIR / (NIR + Red)
# NIR = Band 5, Red = Band 4
vi <- function(img, i, k) {
  bk <- img[[k]]
  bi <- img[[i]]
  vi <- (bk) / (bi + bk)
  return(vi)
}
filenames <- paste0('LC08_L2SP_173049_20131220_20200912_02_T1_SR_B', 1:7, ".tif")
filenames
landsat <- rast(filenames)
landsat
ipvi <- vi(landsat, 5, 4)
ipvi
options(scipen=10000)
colors <- brewer.rdylgn(100)
plot(ipvi, col=colors, font.main = 1, main = "IPVI for Landsat-8 OLI/TIRS C1 image LC08_L2SP_173049_20131220_20200912_02_T1_SR_B: \nWhite Nile and Blue Nile confluence, Sudan (2013)", cex.main=0.9)
minor.tick(nx = 10, ny = 10, tick.ratio = 0.3)
# Plotting the histogram
hist(ipvi, font.main = 1, main = "IPVI values for Landsat-8 OLI/TIRS C1 image \nLC08_L2SP_173049_20131220_20200912_02_T1_SR_B: White Nile and Blue Nile confluence, Sudan (2013)", xlab = "IPVI", ylab= "Frequency",
    col = "yellowgreen", xlim = c(0, 1),  breaks = 50, xaxt = "n")
axis(side=1, at = seq(0.0, 1, 0.1), labels = seq(0.0, 1, 0.1))
minor.tick(nx = 10, ny = 10, tick.ratio = 0.3)
\end{lstlisting}

%%%%%%%%%%%%%%%%%%%%%%%%%%%%%%%%%%%%%%%%%%
\begin{adjustwidth}{-\extralength}{0cm}
%\printendnotes[custom] % Un-comment to print a list of endnotes

\reftitle{References}

%=====================================
% References, variant A: external bibliography
%=====================================
\bibliography{Sudan.bib}
%\nocite{*}

%%%%%%%%%%%%%%%%%%%%%%%%%%%%%%%%%%%%%%%%%%
\end{adjustwidth}
\end{document}
